% Определения терминов (\termsanddefenitions подтягивает их по \newglossaryentry).
\newglossaryentry{silver_labels}{
    name={Серебряная разметка},
    description={набор пар «товар --- целевое значение атрибута», полученный автоматически из открытых источников данных (теги Open Food Facts, OFF) и используемый для обучения моделей. Содержит шум и неполноту; противопоставляется эталонной разметке}
}
\newglossaryentry{gold_labels}{
    name={Эталонная разметка},
    description={пары «товар --- целевое значение», прошедшие верификацию согласованным мнением нескольких больших языковых моделей и (или) слепой разметкой более сильной модели; используется как опорная истина при оценке каскада}
}
\newglossaryentry{weak_supervision}{
    name={Обучение со слабой разметкой (англ.\ weak supervision)},
    description={режим обучения с использованием неполной или зашумлённой разметки, формируемой автоматически из доступных источников вместо ручной разметки экспертом}
}
\newglossaryentry{cascade_arch}{
    name={Каскадная архитектура},
    description={последовательность алгоритмических слоёв, каждый из которых обрабатывает только те ячейки «товар, атрибут», которые не были закрыты предыдущими слоями выше калиброванного порога уверенности}
}
\newglossaryentry{router_term}{
    name={Маршрутизатор (категорий)},
    description={предкаскадный классификатор-распределитель (Слой 0): по партнёр-доступным полям определяет товарную категорию либо относит товар к классу вне области охвата (OOD), после чего запускается схема каскада соответствующей категории}
}
\newglossaryentry{conf_filter}{
    name={Пошаговая фильтрация по уверенности},
    description={формализация поведения каскада как последовательного конъюнктивного условия на прохождение калиброванных порогов уверенности на каждом слое каскада}
}
\newglossaryentry{coverage}{
    name={Покрытие},
    description={доля ячеек «товар, атрибут» в тестовой выборке, для которых каскад дал предсказание выше порога уверенности на слоях 1--3 (без обращения к запасному слою)}
}
\newglossaryentry{covered_acc}{
    name={Точность на покрытом срезе},
    description={доля верных предсказаний среди тех ячеек, для которых каскад дал ответ без обращения к запасному слою}
}
\newglossaryentry{end2end_acc}{
    name={Сквозная точность},
    description={точность системы на полной тестовой выборке, при которой отказы каскада засчитываются как неверные ответы (либо заменяются ответом запасного слоя при включённом Слое 4)}
}
\newglossaryentry{code_disjoint}{
    name={Разбиение без пересечения товарных кодов},
    description={деление выборки на обучающую и тестовую части так, что ни один товарный код (баркод) не попадает одновременно в обе части; обеспечивает оценку качества системы на новых товарных позициях}
}
\newglossaryentry{audit_loop}{
    name={Цикл «аудит --- правки --- переобучение --- измерение»},
    description={методология итеративного улучшения системы: выявленные при аудите дефекты эталонной разметки устраняются правкой экстрактора признаков, после чего модели переобучаются, а прирост точности измеряется на той же контрольной выборке}
}
